\documentclass[10pt,a4paper]{amsart}
\usepackage[margin=19mm]{geometry}
\usepackage{amsmath,amssymb,mathtools}
\usepackage[scr=rsfs,cal=euler]{mathalfa}
\usepackage{xfrac}
\usepackage[unicode,hidelinks,bookmarks=false]{hyperref}
\newcommand{\ie}{i.e.\ }
\newcommand{\cf}{cf.\ }
\newcommand{\resp}{respectively\ }
\newcommand{\Subsection}{\subsection}
\providecommand{\nobreakdash}{\nobreak\hbox{-}}
\setlength{\emergencystretch}{2em}
\multlinegap0pt
\usepackage{amssymb}
\usepackage{mathtools}
\newtheorem{lemma}{Lemma}
\newtheorem{assumption}{Assumption}
\newcommand{\C}{\mathbb C}
\newcommand{\E}{\mathbb E}
\newcommand{\Prob}{\mathbb P}
\newcommand{\logm}{\log^-}
\title{Convergence for Small-Order Derivatives of Random Polynomials with Independent Roots}
\author{Hongkai Zhu}
\date{Source: arXiv:2609.15984v1 (CC BY 4.0)}
\begin{document}
\maketitle

\noindent\emph{Source and licence.} Convergence for Small-Order Derivatives of Random Polynomials with Independent Roots. Version arXiv:2609.15984v1 (CC BY 4.0), distributed by arXiv under the Creative Commons Attribution 4.0 International licence, http://creativecommons.org/licenses/by/4.0/, as recorded in the arXiv licence metadata for this exact version and shown on the abstract page https://arxiv.org/abs/2609.15984v1. Full licence notice: \texttt{licenses/arxiv-2609.15984-CC-BY-4.0.txt}.\par\smallskip

\noindent\textit{Editorial scope.} Counted unit: the article's lemma lem:summable (source0.tex lines 465--469). The article's complete original proof of it is delivered with it (source0.tex lines 471--571, one \texttt{proof} environment of 101 lines). No mathematics is written, repaired or re-derived: the statement and the proof are the article's own lines, in the article's own notation, and no retained line is substituted or re-flowed.  Retained uncounted as the source-local context the proof needs: lines 345--360; lines 431--433.  Nothing was omitted from any retained span, so the package carries no omission record.  No retained line contains a citation, so the wrapper carries no bibliography and no bibliography entry.  No retained line contains a figure environment, an image-inclusion command or drawing code, so no image is delivered and the package declares no asset dependency.  Editorial formatting only, in addition: an AMS \texttt{amsart} preamble that restates the article's own declarations and macros used by the retained lines, namely the environments \texttt{lemma}, \texttt{assumption} on separate counters, together with the standard packages those lines need.  The retained environments print in this document as Lemma 1, Assumption 1 in order of appearance, whereas the article prints the same items with the numbering of its own full document.  The complete native source of the article as distributed in the arXiv e-print for this exact version is bundled unchanged as \texttt{sources/210469/source0.tex}.
\begin{lemma}[Polynomial postselection]\label{lem:postselection}
Let $K\subset\C$ be compact, fix $L<\infty$ and $\theta>0$, and let
$r_N=o(N)$. For every $c>0$, there exists $\rho>0$ depending only on $c, K, L$ with the following property.

Let $\nu$ be any Borel probability measure supported on $K$ such that $L_\nu=\sup_{w\in\C} \int_\C\logm|t-w|\,d\nu(t)\le L$.
For each positive integer $N$, let $M_N$ and $m_N$ satisfy $M_N\le\rho N$ and $m_N\ge\theta N$, let $V_N\subset\C[t]_{\le M_N}:=\{v\in \C[t]:\deg v\le M_N\}$ be a complex linear space with $\dim V_N\le r_N$, and let
$T_1,\ldots,T_{m_N}$ be i.i.d. with law $\nu$. Then
\[\Prob^*\!\left( \exists v\in V_N:\ v(0)=1,\quad |v(T_i)|\le e^{-cN}
     \text{ for every }1\le i\le m_N\right)
 \le e^{-c_1m_N}\]
for all sufficiently large $N$ where $c_1>0$, and one may take $c_1=1/64$.

The choice of $\rho$ and the estimate are uniform over all measures $\nu$
and spaces $V_N$ satisfying the preceding assumptions. The large-$N$
threshold may depend on the fixed parameters and on the sequence $(r_N)$.
\end{lemma}

\begin{assumption}\label{assump}
    Let $X_1,\ldots, X_n$ be i.i.d. with law $\mu$. Let $k_n=o(n)$ be deterministic and assume there exists $a\in (0,1]$ and a compactly supported probability measure $\lambda$ satisfying $a\lambda\le \mu$ and $L_\lambda<\infty$.
\end{assumption} %All probabilities and expectations below refer to these i.i.d. roots. Such a component exists whenever $\mu$ is not polar carries, by proposition xx. (Also below)

\begin{lemma}\label{lem:summable}
    Under Assumption~\ref{assump}, for each fixed bounded disk $D$, $\varepsilon>0$, and $0<\delta<1/2$, there exists $T_0$ such that for every fixed $T\ge T_0$, there is $c_T>0$ for which 
    \[\E|\{z\in G_T:B_n(z)\}|\le |D|e^{-c_T n}\]
    for all sufficiently large $n$. In particular, these expectations are summable in $n$.
\end{lemma}

\begin{proof}
    Fix $z\in G_T$. There is nothing to prove when $z$ is an atom of $\mu$, otherwise $P_n(z)\neq 0$ almost surely. Also if $k_n=0$, then $B_n(z)=\emptyset$ because $R_{n,0}=1$. We therefore assume $z$ is non-atom and $k_n\ge 1$. Since $z\in G_T$,
    \[\lambda(\C\setminus B(z,e^{-T}))\ge\frac{1}{2},
    \qquad \mu(B(z,e^{-T}))\le e^{-T}.\]

    Set 
    \[S_z:=\C\setminus B(z,e^{-T}),\qquad p_z:=a\lambda(S_z).\]
    Since $a\lambda\le \mu$, we have $a\lambda\mid_{S_z}\le \mu$. Moreover, because $z\in G_T$, 
    \[\lambda(S_z)=1-\lambda(B(z,e^{-T}))\ge \frac{1}{2},\]
    and hence $p_z\ge a/2$. Define the probability measure 
    \[\lambda_z^{\mathrm{out}}:=\frac{\lambda\mid_{S_z}}{\lambda(S_z)}.\]
    If $p_z<1$, define 
    \[\rho_z:=\frac{\mu-a\lambda\mid_{S_z}}{1-p_z}.\]
    If $p_z=1$, let $\rho_z$ be any probability measure. In either case,
    \[\mu=p_z\lambda_z^{\mathrm{out}}+(1-p_z)\rho_z.\]

    We may therefore generate the roots independently as follows. For each $j$, let $\eta_j$ be Bernoulli with $\Prob(\eta_j=1)=p_z$. Conditional on $\eta_j=1$, sample $X_j$ with law $\lambda_z^{\mathrm{out}}$; conditional on $\eta_j=0$, sample $X_j$ with law $\rho_z$. The resulting variables $X_1,\ldots, X_n$ are still i.i.d. with law $\mu$. We call $X_j$ eligible when $\eta_j=1$. Thus the eligibility labels are independent, each root is eligible with probability $p_z\ge a/2$. Condition on all the labels, the eligible roots are independent with common law $\lambda_z^{\mathrm{out}}$.
    
    Let $m=\lceil \theta n\rceil$ where $\theta>0$ will be chosen later. Since the number of eligible roots has distribution $\operatorname{Bin}(n,p_z)$ with $p_z\ge a/2$, if $\theta<a/4$, Chernoff bound gives 
    \[\Prob(\text{there are fewer than }m\text{ eligible roots})=\Prob(\#\{1\le j\le n:\eta_j=1\}<m)\le e^{-an/16}.\]
    Since also $\mu(B(z,e^{-T}))\le e^{-T}$, Chernoff bound again gives that except on an event of probability at most $e^{-e^{-T}n/3}$ there are at most $2e^{-T}n$ roots in $B(z,e^{-T})$, i.e.
    \[\Prob(\#\{j:X_j\in B(z,e^{-T})\}>2e^{-T}n)\le e^{-e^{-T}n/3}.\]
    Call the roots in this ball the \emph{near roots}.
    
    On the intersection of these two good events, retain the first $m$ eligible roots and condition on all mixture labels and all unretained roots. Since every retained root is eligible, it lies outside $B(z,e^{-T})$; therefore the near-root count is determined entirely by the unretained roots and is part of the conditioned background. Thus, conditionally on this background, the retained roots $Y_1,\ldots, Y_m$ remain i.i.d. with law $\lambda_z^{\mathrm{out}}$. Therefore the translates $Y_i-z$ are conditionally i.i.d. with the corresponding translated law $\lambda_z^{\mathrm{out}}$. Define
    \[L:=\max\left\{1,\sup_{x\in \operatorname{supp} \lambda,z\in \overline{D}}|x-z|\right\}.\]
    Their law is supported on $\overline{B(0,L)}$ and satisfies, uniformly in $z\in D$,
    \[\sup_{w\in\C}\frac{1}{\lambda(S_z)}\int_{\C\setminus B(z,e^{-T})}\log^-|x-z-w|\,d\lambda(x)
    \le 2\sup_{u\in\C}\int_\C\log^-|x-u|\,d\lambda(x) 
    \le 2L_\lambda.\]
    Indeed, translation only changes the point in the supremum, restriction can only decrease the integral.

    Put
    \[F(t):=\prod_{j=1}^n\left(1+\frac{t}{z-X_j}\right)=\sum_{r=0}^n \binom{n}{r}R_{n,r}(z)t^r.\]
    Therefore on $B_n(z)$, the coefficient of $t^r$ satisfies $|\binom{n}{r}R_{n,r}|\le \binom{n}{r}e^{-\varepsilon n}$ for $k_n\le r\le k_n+\lfloor \delta n\rfloor$.
    We now invert only the factors corresponding to unretained roots outside $B(z,e^{-T})$. Define 
    \[Q(t):=\prod_{\substack{j\text{ unretained}\\
                     X_j\notin B(z,e^{-T})}}
    \left(1+\frac{t}{z-X_j}\right)^{-1}=\sum_{\ell\ge 0}b_\ell t^\ell.\]
    Each inverted factor has constant term $1$, so its formal inverse is well-defined. The series also converges absolutely for $|t|<e^{-T}$, but below we use only its coefficients and finite truncations. Because $Q$ depends only on the unretained roots, it is fixed once the background is conditioned on.

    Cancellation in the formal product give
    \[Q(t)F(t)=
    \prod_{i=1}^{m}\left(1+\frac{t}{z-Y_i}\right)\prod_{\substack{j\ \mathrm{unretained}\\X_j\in B(z,e^{-T})}}\left(1+\frac{t}{z-X_j}\right).\]
    Thus $QF$ is a polynomial. On the good counting event,
    \[\deg(QF)\leq M:=m+\lceil 2e^{-T}n\rceil.\]
    In particular, evaluating this polynomial gives
    \[(QF)(Y_i-z)=0,\qquad 1\leq i\leq m.\]
    These evaluations do not require convergence of the power series for $Q$ at $Y_i-z$.
    
    Every factor inverted in $Q$ corresponds to a root outside $B(z,e^{-T})$ and therefore satisfies $|\frac{1}{z-X_j}| \le e^T$. Expanding the reciprocal factors as geometric series gives
    \[|b_\ell| \le \binom{n+\ell-1}{\ell}e^{T\ell}, \qquad \ell\ge0.\]
    
    Write $(\cdot)_{\le M}$ for truncation through degree $M$, and define the background-determined linear space
    \[V:=\left\{(Qh)_{\le M}:h\in\C[t]_{\le k_n-1}\right\}
    \subset \C[t]_{\le M},\qquad\dim V\le k_n=o(n).\]
    Define 
    \[v:=\left(Q\sum_{s=0}^{k_n-1}\binom{n}{s}R_{n,s}(z)t^s\right)_{\le M}.\]
    Then $v\in V$ and $v(0)=1$.
    Although $V$ is fixed by the conditioned background, the particular polynomial $v$ may depend on the retained roots. This is precisely the postselection allowed by Lemma~\ref{lem:postselection}.
    
    Since $QF$ has degree at most $M$ on the good event, the definition of $v$ gives
    \[(QF-v)(t)
    =\left(Q(t)\sum_{s=k_n}^{n}\binom nsR_{n,s}(z)t^s\right)_{\le M} 
    =\sum_{r=k_n}^{M}\left(\sum_{s=k_n}^{r}b_{r-s}\binom nsR_{n,s}(z)\right)t^r.\]
    Suppose that $M\le k_n+\lfloor\delta n\rfloor$, as will follow from the parameter choice below. On $B_n(z)$, for every $k_n\le r\le M$, which is nonempty since $M\ge m\ge\theta n$ and $k_n=o(n)$,
    \[\left|\sum_{s=k_n}^{r}b_{r-s}\binom nsR_{n,s}(z)\right|\le(M+1)e^{-\varepsilon n}\max_{\substack{k_n\le s\le M\\0\le\ell\le M}}\binom ns\binom{n+\ell-1}{\ell}e^{T\ell}.\]
    Since both $QF$ and $v$ have degree at most $M$ and $(QF)(Y_i-z)=0$, on the intersection of $B_n(z)$ with the good event,
    \[|v(Y_i-z)|\le(M+1)^2L^M e^{-\varepsilon n}\max_{\substack{k_n\le s\le M\\0\le\ell\le M}}\binom ns\binom{n+\ell-1}{\ell}e^{T\ell},
    \qquad 1\le i\le m.\]
    Set $\alpha:=\theta+3e^{-T}$. For all sufficiently large $n$, $M/n\le\alpha$. If $\alpha<1/2$, the elementary binomial estimate $\binom Nr\le\left(\frac{eN}{r}\right)^r$ gives, for $0\le s,\ell\le M$,
    \[\binom ns\binom{n+\ell-1}{\ell}\le\exp\left\{Cn\alpha\log\frac e\alpha\right\}\]
    for some absolute constant $C>0$. And since $L^Me^{T\ell}\le\exp\{n\alpha(T+\log L)\}$, on $B_n(z)$ and the good event,
    \[|v(Y_i-z)|\le\exp\left(-n\left(\varepsilon-C\alpha\log\frac{e}{\alpha}-(T+\log L)\alpha-o(1)\right)\right),
    \qquad 1\le i\le m.\]
    Here the $o(1)$ absorbs the factor $(M+1)^2$ for all $z$.
    
    Let $\rho_*>0$ be supplied by Lemma~\ref{lem:postselection} with $c=\varepsilon/2$, compact support $\overline{B(0,L)}$, and potential bound $2L_\lambda$. Since $e^{-T}\log(e/e^{-T})\to 0$ and $Te^{-T}\to 0$, we may choose $T_0$ sufficiently large that, for every $T\ge T_0$,
    \[3e^{-T}<\min\{\rho_*,\delta/2,1/2\}
    \qquad\text{and}\qquad
    C(3e^{-T})\log\frac{e}{3e^{-T}}+(T+\log L)3e^{-T}<\frac{\varepsilon}{4}.\]

    Fix $T\ge T_0$. Since $x\log(e/x)\to 0$ as $x\to 0$ and $\alpha=\theta+3e^{-T}\to3e^{-T}$ as $\theta\to 0$, we may choose $\theta>0$ sufficiently small that
    \[\theta<\frac a4,
    \qquad
    \alpha<\min\{\rho_*,\delta/2,1/2\},
    \qquad
    C\alpha\log\frac e\alpha+(T+\log L)\alpha<\frac{\varepsilon}{3}.\]
    For all sufficiently large $n$, these choices imply $M\le\rho_*n$, $M\le k_n+\lfloor\delta n\rfloor$. They also give
    \[|v(Y_i-z)|\le e^{-\varepsilon n/2}, \qquad 1\le i\le m\] on $B_n(z)$ and the good event.
    
    Now fix any background satisfying the two good conditions. Under this conditioning, $V$ is fixed and $Y_1-z,\ldots,Y_m-z$ are i.i.d. with the translated law described above. Moreover, $V\subset\C[t]_{\le M}$, $\dim V\le k_n=o(n)$, $M\le\rho_*n$, $m\ge\theta n$. On the measurable section of $B_n(z)$, the preceding construction supplies a polynomial $v\in V$ satisfying $v(0)=1$ and $|v(Y_i-z)|\le e^{-\varepsilon n/2}$ for every $1\le i\le m$. 

    Therefore this measurable section is contained in the existential event of Lemma~\ref{lem:postselection}. Apply that lemma under the conditional product law of $Y_1-z,\ldots,Y_m-z$. Under this law, the probability of the measurable section is at most the outer probability of the containing event, and hence at most $e^{-c_1m}$. The estimate and the large-$n$ threshold are uniform over the conditioned backgrounds and $z\in G^T$. Since the two good conditions are determined by the background, averaging these ordinary conditional probabilities and adding the two Chernoff exceptions gives
    \[\Prob(B_n(z))
    \le e^{-c_1m}+e^{-an/16}+e^{-e^{-T}n/3}
    \le e^{-c_Tn}\]
    for some $c_T>0$ and all sufficiently large $n$. Finally Tonelli's theorem gives
    \[\E\bigl|\{z\in G_T:B_n(z)\}\bigr|=\int_{G_T}\Prob(B_n(z))\,d^2z \le|D|e^{-c_Tn}\]
    which completes the proof.
\end{proof}

\end{document}
